\section{LLM-planned executed cover}\label{supp:cover-study}
The experiment measures traffic recognition, requested-task completion and full business-state integrity on the same executions. It uses six synthetic ticket actions and six policies: no cover, proxy discard, read cover, hidden isolated execution, live writes and visible isolated execution. All policies share fixtures and frozen random cover draws. The earlier reference-program replay remains a separate artifact; its fixed-position and one-to-four-decoy sweeps are not LLM results.

\subsection{Population, planner and cover semantics}
Each of two fixed-weight models, Qwen3:8B and Llama3.1:8B, runs 300 fixtures per condition: 120 training, 60 validation and 120 test, balanced across six requested actions. This gives 3,600 episodes, including 1,440 test executions of 120 unique paired test fixtures. Request phrasings recur across splits; this is fixture holdout, not unseen-template transfer. The seed pairs fixtures and decoy schedules, and condition order rotates. Two workers collect each model cohort on one host.

LangGraph alternates model decisions with actual tool calls. The controller fetches the ticket; required checks constrain action admissibility without naming the correct action or object. The model chooses check order, requested business operation, order identifier, recovery and finish, up to ten decision steps. Temperature is zero, context 4,096 and maximum generated decision length 96 tokens; the fixture supplies the model seed. Fixed manifests and weight blobs match before and after collection.

The fixture contains three orders, customer address, a compensation quota, escalation status and a ledger. Its intended reference is the initial state after exactly the requested legitimate action. Decoys deliberately include operations capable of collateral changes. A frozen schedule places one or two cover calls after the ticket or early honest calls; extra later planner calls do not add cover. Each cover cell triggers all 182 planned test decoys.

Proxy discard computes the response on a disposable copy, retaining no decoy state. Read cover queries state without changing it. Hidden isolated execution applies the decoy on a persistent clone; live write applies it to the primary fixture. Visible isolated execution uses the same clone and schedule, forwarding clearly marked gateway-cover requests and responses to the planner. These objects are not fabricated planner tool calls. In the other conditions, cover responses are hidden from the planner.

\subsection{Transport, observer and endpoints}
Seven local MCP Streamable HTTP destinations use actual TLS~1.3/HTTP/2 sockets and a common \path{/mcp} path. Request bodies use a 4,096-byte padding quantum and responses 16,384 bytes. All seven sessions initialize before each episode and are reused within it. Total TLS cost includes those initializations; call-byte cost and attack features exclude them. The observer receives destination and bidirectional ciphertext byte counts in supplied episode and logical-call windows. It receives no plaintext, model traffic, physical timing, arm label, fixture identifier or private stream assignment. The original replay's variable HTTP-path artifact therefore does not carry over.

Recognition targets the requested action, retaining wrong or failed executions. Fixed features encode the first sixteen call destinations and ciphertext byte counts, per-destination totals and total calls/bytes. LR uses training-only standardization, $C=1$ and a 3,000-iteration cap. RF uses 300 trees with fixed seed. Validation balanced accuracy selects LR or RF within each model/condition, breaking ties toward LR. Both candidates, the no-cover-selected frozen attack and last/penultimate destination rules remain reported (Table~\ref{tab:m3-llm-recognition}).

Goal completion requires the requested action/object, satisfied prerequisites and the target postcondition. State integrity compares the full business state, including non-target orders, shared fields and the ledger, against the single intended-action reference. The predicates differ: a reference-matching state can still fail the stricter action/object-history goal predicate. Conversely, a completed goal can leave collateral changes. Table~\ref{tab:m3-llm-utility} reports both and their completed-but-incorrect-state intersection.

\begin{table*}[t]
\centering\small
\setlength{\tabcolsep}{4pt}
\caption{LLM-planned executed cover on 120 paired held-out fixtures per cell. The primary LR/RF attacker is selected on validation data. All attack columns count correct predictions out of 120; chance is 20/120. The interval conditions on the fitted attacker and synthetic fixtures.}
\label{tab:m3-llm-recognition}
\begin{tabular}{llrrrrrl}
\toprule
Model & Cover & Frozen & LR & RF & Last & Penultimate & Selected BA [95\% CI] \\
\midrule
Qwen3:8B & No cover & 118 & 118 & 117 & 116 & 14 & 0.983 [0.958, 1.000] (LR) \\
Qwen3:8B & Proxy discard & 44 & 63 & 62 & 77 & 43 & 0.517 [0.425, 0.608] (RF) \\
Qwen3:8B & Read cover & 42 & 59 & 62 & 77 & 42 & 0.517 [0.425, 0.608] (RF) \\
Qwen3:8B & Isolated, hidden & 43 & 62 & 61 & 77 & 42 & 0.508 [0.417, 0.600] (RF) \\
Qwen3:8B & Live write & 42 & 60 & 62 & 75 & 44 & 0.517 [0.425, 0.608] (RF) \\
Qwen3:8B & Isolated, visible & 40 & 63 & 63 & 77 & 42 & 0.525 [0.433, 0.617] (RF) \\
Llama3.1:8B & No cover & 113 & 113 & 110 & 79 & 64 & 0.942 [0.900, 0.975] (LR) \\
Llama3.1:8B & Proxy discard & 54 & 81 & 79 & 67 & 61 & 0.658 [0.575, 0.733] (RF) \\
Llama3.1:8B & Read cover & 57 & 78 & 72 & 67 & 62 & 0.600 [0.517, 0.683] (RF) \\
Llama3.1:8B & Isolated, hidden & 56 & 81 & 82 & 67 & 60 & 0.683 [0.600, 0.758] (RF) \\
Llama3.1:8B & Live write & 58 & 80 & 75 & 65 & 59 & 0.667 [0.583, 0.750] (LR) \\
Llama3.1:8B & Isolated, visible & 58 & 79 & 81 & 68 & 64 & 0.675 [0.592, 0.750] (RF) \\
\bottomrule
\end{tabular}
\end{table*}

\begin{table*}[t]
\centering\small
\setlength{\tabcolsep}{4pt}
\caption{Task outcomes and measured costs on the same 120 held-out fixtures. State integrity compares the entire business state with the single intended-action reference. Byte ratios are means of paired ratios; total TLS bytes include all seven MCP initializations, while call bytes exclude initialization.}
\label{tab:m3-llm-utility}
\begin{tabular}{llrrrrrr}
\toprule
Model & Cover & Goal & State & Goal + state error & Call bytes $\times$ & Total bytes $\times$ & Model calls \\
\midrule
Qwen3:8B & No cover & 114/120 & 116/120 & 0/120 & 1.00 & 1.00 & 4.06 \\
Qwen3:8B & Proxy discard & 114/120 & 116/120 & 0/120 & 1.34 & 1.15 & 4.06 \\
Qwen3:8B & Read cover & 114/120 & 115/120 & 0/120 & 1.34 & 1.15 & 4.07 \\
Qwen3:8B & Isolated, hidden & 114/120 & 115/120 & 0/120 & 1.34 & 1.15 & 4.07 \\
Qwen3:8B & Live write & 102/120 & 7/120 & 95/120 & 1.34 & 1.15 & 4.08 \\
Qwen3:8B & Isolated, visible & 113/120 & 114/120 & 1/120 & 1.35 & 1.15 & 4.08 \\
Llama3.1:8B & No cover & 82/120 & 69/120 & 13/120 & 1.00 & 1.00 & 5.23 \\
Llama3.1:8B & Proxy discard & 80/120 & 68/120 & 13/120 & 1.29 & 1.14 & 5.31 \\
Llama3.1:8B & Read cover & 84/120 & 71/120 & 14/120 & 1.28 & 1.13 & 5.25 \\
Llama3.1:8B & Isolated, hidden & 83/120 & 70/120 & 13/120 & 1.28 & 1.13 & 5.25 \\
Llama3.1:8B & Live write & 76/120 & 3/120 & 73/120 & 1.29 & 1.14 & 5.35 \\
Llama3.1:8B & Isolated, visible & 83/120 & 73/120 & 11/120 & 1.29 & 1.14 & 5.33 \\
\bottomrule
\end{tabular}
\end{table*}


\subsection{Paired inference and uncertainty}
The analysis fixes five cover-minus-none contrasts plus visible-minus-hidden for each model. Each endpoint---recognition, completion and integrity---has its own twelve-comparison Holm family. Two-sided exact discordant-pair tests use paired fixture outcomes. Intervals use 10,000 paired fixture bootstrap draws stratified by requested action and condition on the fitted attackers and designed workload. They do not represent template, fit-selection or deployment uncertainty. Table~\ref{tab:m3-llm-paired} retains every prespecified comparison, including null contrasts.

All ten primary recognition contrasts versus no cover survive correction; the largest adjusted $p$ is $3.08\times10^{-8}$. Live-write integrity reductions survive for both models ($3.70\times10^{-32}$ and $2.98\times10^{-19}$), while completion survives only for Qwen ($.02197$, versus $1$ for Llama). Every primary visible-minus-hidden contrast has adjusted $p=1$. These tests do not establish equivalence among cover policies or a privacy upper bound. In particular, Qwen's last-destination rule exceeds the validation-selected learner under cover.

\begin{table*}[t]\centering\small
\caption{Goal completion: every prespecified paired contrast. Differences and conditional intervals are percentage points. This endpoint has one twelve-comparison Holm family.}
\label{tab:m3-llm-paired}
\setlength{\tabcolsep}{5pt}
\begin{tabularx}{\textwidth}{@{}lYrrrr@{}}\toprule
Model & Condition minus reference & Difference & Conditional 95\% CI & Raw $p$ & Holm $p$ \\\midrule
Qwen3:8B & Proxy -- No cover & 0.00 & [0.00, 0.00] & 1.000 & 1.000 \\
Qwen3:8B & Read -- No cover & 0.00 & [0.00, 0.00] & 1.000 & 1.000 \\
Qwen3:8B & Hidden clone -- No cover & 0.00 & [0.00, 0.00] & 1.000 & 1.000 \\
Qwen3:8B & Live write -- No cover & -10.00 & [-15.00, -5.00] & 0.002 & 0.022 \\
Qwen3:8B & Visible clone -- No cover & -0.83 & [-4.17, 2.50] & 1.000 & 1.000 \\
Qwen3:8B & Visible clone -- Hidden clone & -0.83 & [-4.17, 2.50] & 1.000 & 1.000 \\
\midrule
Llama3.1:8B & Proxy -- No cover & -1.67 & [-5.00, 1.67] & 0.625 & 1.000 \\
Llama3.1:8B & Read -- No cover & 1.67 & [-2.50, 5.83] & 0.688 & 1.000 \\
Llama3.1:8B & Hidden clone -- No cover & 0.83 & [0.00, 2.50] & 1.000 & 1.000 \\
Llama3.1:8B & Live write -- No cover & -5.00 & [-10.83, 0.83] & 0.210 & 1.000 \\
Llama3.1:8B & Visible clone -- No cover & 0.83 & [-6.67, 8.33] & 1.000 & 1.000 \\
Llama3.1:8B & Visible clone -- Hidden clone & 0.00 & [-7.50, 7.50] & 1.000 & 1.000 \\
\bottomrule\end{tabularx}
\end{table*}
\begin{table*}[t]\centering\small
\caption{Requested-action recognition: every prespecified paired contrast. Differences and conditional intervals are percentage points. This endpoint has one twelve-comparison Holm family.}
\label{tab:m3-llm-paired-privacy}
\setlength{\tabcolsep}{5pt}
\begin{tabularx}{\textwidth}{@{}lYrrrr@{}}\toprule
Model & Condition minus reference & Difference & Conditional 95\% CI & Raw $p$ & Holm $p$ \\\midrule
Qwen3:8B & Proxy -- No cover & -46.67 & [-55.83, -37.50] & $2.78\times10^{-17}$ & $3.05\times10^{-16}$ \\
Qwen3:8B & Read -- No cover & -46.67 & [-55.00, -37.50] & $2.78\times10^{-17}$ & $3.05\times10^{-16}$ \\
Qwen3:8B & Hidden clone -- No cover & -47.50 & [-56.67, -38.33] & $1.39\times10^{-17}$ & $1.67\times10^{-16}$ \\
Qwen3:8B & Live write -- No cover & -46.67 & [-55.83, -37.50] & $2.78\times10^{-17}$ & $3.05\times10^{-16}$ \\
Qwen3:8B & Visible clone -- No cover & -45.83 & [-55.00, -36.67] & $8.05\times10^{-16}$ & $6.44\times10^{-15}$ \\
Qwen3:8B & Visible clone -- Hidden clone & 1.67 & [-5.00, 8.33] & 0.815 & 1.000 \\
\midrule
Llama3.1:8B & Proxy -- No cover & -28.33 & [-35.83, -20.83] & $1.08\times10^{-9}$ & $5.59\times10^{-9}$ \\
Llama3.1:8B & Read -- No cover & -34.17 & [-42.50, -25.83] & $9.09\times10^{-13}$ & $6.37\times10^{-12}$ \\
Llama3.1:8B & Hidden clone -- No cover & -25.83 & [-34.17, -18.33] & $9.31\times10^{-10}$ & $5.59\times10^{-9}$ \\
Llama3.1:8B & Live write -- No cover & -27.50 & [-35.83, -19.17] & $1.02\times10^{-8}$ & $3.07\times10^{-8}$ \\
Llama3.1:8B & Visible clone -- No cover & -26.67 & [-35.00, -19.17] & $4.07\times10^{-9}$ & $1.63\times10^{-8}$ \\
Llama3.1:8B & Visible clone -- Hidden clone & -0.83 & [-8.33, 6.67] & 1.000 & 1.000 \\
\bottomrule\end{tabularx}
\end{table*}
\begin{table*}[t]\centering\small
\caption{Business-state integrity: every prespecified paired contrast. Differences and conditional intervals are percentage points. This endpoint has one twelve-comparison Holm family.}
\label{tab:m3-llm-paired-state_integrity}
\setlength{\tabcolsep}{5pt}
\begin{tabularx}{\textwidth}{@{}lYrrrr@{}}\toprule
Model & Condition minus reference & Difference & Conditional 95\% CI & Raw $p$ & Holm $p$ \\\midrule
Qwen3:8B & Proxy -- No cover & 0.00 & [0.00, 0.00] & 1.000 & 1.000 \\
Qwen3:8B & Read -- No cover & -0.83 & [-2.50, 0.00] & 1.000 & 1.000 \\
Qwen3:8B & Hidden clone -- No cover & -0.83 & [-2.50, 0.00] & 1.000 & 1.000 \\
Qwen3:8B & Live write -- No cover & -90.83 & [-95.00, -85.83] & $3.08\times10^{-33}$ & $3.70\times10^{-32}$ \\
Qwen3:8B & Visible clone -- No cover & -1.67 & [-5.83, 1.67] & 0.688 & 1.000 \\
Qwen3:8B & Visible clone -- Hidden clone & -0.83 & [-4.17, 2.50] & 1.000 & 1.000 \\
\midrule
Llama3.1:8B & Proxy -- No cover & -0.83 & [-4.17, 2.50] & 1.000 & 1.000 \\
Llama3.1:8B & Read -- No cover & 1.67 & [-2.50, 6.67] & 0.727 & 1.000 \\
Llama3.1:8B & Hidden clone -- No cover & 0.83 & [-1.67, 3.33] & 1.000 & 1.000 \\
Llama3.1:8B & Live write -- No cover & -55.00 & [-62.50, -47.50] & $2.71\times10^{-20}$ & $2.98\times10^{-19}$ \\
Llama3.1:8B & Visible clone -- No cover & 3.33 & [-4.17, 10.83] & 0.523 & 1.000 \\
Llama3.1:8B & Visible clone -- Hidden clone & 2.50 & [-5.00, 10.00] & 0.664 & 1.000 \\
\bottomrule\end{tabularx}
\end{table*}


\subsection{Planner paths, costs and limitations}
Hidden and visible isolated runs have identical cover responses in all 120 paired fixtures per model. Their honest action sequences differ in 16 Qwen and 65 Llama pairs. Hidden-arm comparisons also vary: proxy versus hidden-isolated differs in one Qwen and fifteen Llama pairs; read versus hidden-isolated differs in zero and fifteen. Hence individual path differences do not by themselves identify a visibility effect. Mean prompt tokens rise from 1,668.34 to 1,969.49 for Qwen and from 2,251.48 to 2,675.95 for Llama when responses are visible, increases of 18.1\% and 18.9\%.

Cover call-byte ratios are 1.280--1.346; total TLS ratios are 1.131--1.149 because session initialization is charged to both conditions. Mean local wall times are 8.14--9.94 seconds per test episode, including setup, journaling and two-worker contention. These are local resource measurements. Test outputs contain no invalid decision JSON. Llama reaches the step limit once in no-cover, proxy, read and hidden-isolated conditions, twice in live-write and zero times in visible-isolated; Qwen never does. Every terminal case remains included.

All 3,600 formal episodes, 16,601 model request/response chains, 21,163 tool calls and 4,550 cover calls pass separate state, request and transport checks. Formal unresolved requests are zero. One earlier development request remains unknown, was not retried and is excluded from the formal cohort. The release supplies synthetic planner records, fixture states, lengths, predictions and replay scripts; it excludes credentials, model weights, private TLS keys and raw packet captures.

The adverse writes make the consequence measurable rather than estimating its production prevalence. The experiment is neither a reproduction of WRA's research decoys nor a measurement of physical-timing, model-link or WAN privacy. Llama's unprotected completion and state baselines are already imperfect. The positive finding is the observed separation between goal completion and collateral state integrity; visible-response averages remain inconclusive.
